\documentclass[10pt,a4paper]{amsart}
\usepackage{amsmath,amssymb,amsthm,graphicx}
\usepackage{verbatim}
\newenvironment{reference}{\comment}{\endcomment}
%\newenvironment{reference}{}{}
\newenvironment{slogan}{\comment}{\endcomment}
\newenvironment{history}{\comment}{\endcomment}

% For commutative diagrams we use Xy-pic
\usepackage[all]{xy}

% We use 2cell for 2-commutative diagrams.
\xyoption{2cell}
% Theorem environments.
%
\theoremstyle{plain}
\newtheorem{theorem}[subsection]{Theorem}
\newtheorem{proposition}[subsection]{Proposition}
\newtheorem{lemma}[subsection]{Lemma}

\theoremstyle{definition}
\newtheorem{definition}[subsection]{Definition}
\newtheorem{example}[subsection]{Example}
\newtheorem{exercise}[subsection]{Exercise}
\newtheorem{situation}[subsection]{Situation}

\theoremstyle{remark}
\newtheorem{remark}[subsection]{Remark}
\newtheorem{remarks}[subsection]{Remarks}

\numberwithin{equation}{subsection}

% Macros
%
\def\lim{\mathop{\mathrm{lim}}\nolimits}
\def\colim{\mathop{\mathrm{colim}}\nolimits}
\def\Spec{\mathop{\mathrm{Spec}}}
\def\Hom{\mathop{\mathrm{Hom}}\nolimits}
\def\Ext{\mathop{\mathrm{Ext}}\nolimits}
\def\SheafHom{\mathop{\mathcal{H}\!\mathit{om}}\nolimits}
\def\SheafExt{\mathop{\mathcal{E}\!\mathit{xt}}\nolimits}
\def\Sch{\mathit{Sch}}
\def\Mor{\mathop{\mathrm{Mor}}\nolimits}
\def\Ob{\mathop{\mathrm{Ob}}\nolimits}
\def\Sh{\mathop{\mathit{Sh}}\nolimits}
\def\NL{\mathop{N\!L}\nolimits}
\def\CH{\mathop{\mathrm{CH}}\nolimits}
\def\proetale{{pro\text{-}\acute{e}tale}}
\def\etale{{\acute{e}tale}}
\def\QCoh{\mathit{QCoh}}
\def\Ker{\mathop{\mathrm{Ker}}}
\def\Im{\mathop{\mathrm{Im}}}
\def\Coker{\mathop{\mathrm{Coker}}}
\def\Coim{\mathop{\mathrm{Coim}}}

% Boxtimes
%
\DeclareMathSymbol{\boxtimes}{\mathbin}{AMSa}{"02}

%
% Macros for moduli stacks/spaces
%
\def\QCohstack{\mathcal{QC}\!\mathit{oh}}
\def\Cohstack{\mathcal{C}\!\mathit{oh}}
\def\Spacesstack{\mathcal{S}\!\mathit{paces}}
\def\Quotfunctor{\mathrm{Quot}}
\def\Hilbfunctor{\mathrm{Hilb}}
\def\Curvesstack{\mathcal{C}\!\mathit{urves}}
\def\Polarizedstack{\mathcal{P}\!\mathit{olarized}}
\def\Complexesstack{\mathcal{C}\!\mathit{omplexes}}
% \Pic is the operator that assigns to X its picard group, usage \Pic(X)
% \Picardstack_{X/B} denotes the Picard stack of X over B
% \Picardfunctor_{X/B} denotes the Picard functor of X over B
\def\Pic{\mathop{\mathrm{Pic}}\nolimits}
\def\Picardstack{\mathcal{P}\!\mathit{ic}}
\def\Picardfunctor{\mathrm{Pic}}
\def\Deformationcategory{\mathcal{D}\!\mathit{ef}}
\usepackage{cleveref}
\crefname{originallemma}{original lemma}{original lemmas}
\Crefname{originallemma}{Original lemma}{Original lemmas}
\makeatletter
\newlabel{algebra-lemma-complete-henselian}{{04GM}{source}{}{sourcepointer-1}{}}
\newlabel{algebra-lemma-complete-henselian@cref}{{[originallemma][0][]04GM}{[1][0][]source}{}{}{}}
\makeatother
\crefname{originallemma}{original lemma}{original lemmas}
\Crefname{originallemma}{Original lemma}{Original lemmas}
\makeatletter
\newlabel{algebra-lemma-regular-complete-containing-coefficient-field}{{0C0S}{source}{}{sourcepointer-2}{}}
\newlabel{algebra-lemma-regular-complete-containing-coefficient-field@cref}{{[originallemma][0][]0C0S}{[1][0][]source}{}{}{}}
\makeatother
\crefname{originalequation}{original equation}{original equations}
\Crefname{originalequation}{Original equation}{Original equations}
\makeatletter
\newlabel{equation-multicross}{{0C1U}{source}{}{sourcepointer-3}{}}
\newlabel{equation-multicross@cref}{{[originalequation][0][]0C1U}{[1][0][]source}{}{}{}}
\makeatother
\crefname{originaldefinition}{original definition}{original definitions}
\Crefname{originaldefinition}{Original definition}{Original definitions}
\makeatletter
\newlabel{more-algebra-definition-number-of-branches}{{0C26}{source}{}{sourcepointer-4}{}}
\newlabel{more-algebra-definition-number-of-branches@cref}{{[originaldefinition][0][]0C26}{[1][0][]source}{}{}{}}
\makeatother
\crefname{originallemma}{original lemma}{original lemmas}
\Crefname{originallemma}{Original lemma}{Original lemmas}
\makeatletter
\newlabel{more-algebra-lemma-one-dimensional-number-of-branches}{{0C2D}{source}{}{sourcepointer-5}{}}
\newlabel{more-algebra-lemma-one-dimensional-number-of-branches@cref}{{[originallemma][0][]0C2D}{[1][0][]source}{}{}{}}
\makeatother
\crefname{originallemma}{original lemma}{original lemmas}
\Crefname{originallemma}{Original lemma}{Original lemmas}
\makeatletter
\newlabel{varieties-lemma-delta-number-branches-inequality}{{0C43}{source}{}{sourcepointer-6}{}}
\newlabel{varieties-lemma-delta-number-branches-inequality@cref}{{[originallemma][0][]0C43}{[1][0][]source}{}{}{}}
\makeatother
\crefname{originallemma}{original lemma}{original lemmas}
\Crefname{originallemma}{Original lemma}{Original lemmas}
\makeatletter
\newlabel{varieties-lemma-delta-number-branches-inequality-sh}{{0C42}{source}{}{sourcepointer-7}{}}
\newlabel{varieties-lemma-delta-number-branches-inequality-sh@cref}{{[originallemma][0][]0C42}{[1][0][]source}{}{}{}}
\makeatother
\crefname{originallemma}{original lemma}{original lemmas}
\Crefname{originallemma}{Original lemma}{Original lemmas}
\makeatletter
\newlabel{varieties-lemma-delta-same-after-completion}{{0C3W}{source}{}{sourcepointer-8}{}}
\newlabel{varieties-lemma-delta-same-after-completion@cref}{{[originallemma][0][]0C3W}{[1][0][]source}{}{}{}}
\makeatother
\title{Delta-invariant lower bound versus branch count, with equality giving the completed multicross normal form}
\author{The Stacks Project Authors}
\date{stacks-project:curves.tex:local-snapshot-20261012; GNU FDL 1.2 or later}
\begin{document}
\maketitle
\noindent\textit{Editorial note.} All delivered mathematical text is the author's own, copied from the pinned Stacks Project chapter file curves.tex; no statement or proof step was rewritten, reordered or removed.\par\medskip
\noindent\textit{Editorial note.} The counted proof is the author's complete argument: the branch identification, the inequality, the completion reduction with invariance, the strict henselianness, the wedge-of-DVRs structure, and the k[[t]] normalization are all retained verbatim.\par\medskip
\noindent\textit{Editorial note.} The label lemma-multicross-algebra names the counted statement, and the proof environment immediately following it in the source is delivered verbatim as its proof.\par\medskip
\noindent\textit{Editorial note.} The counted statement is carried by the author's own lemma environment in the chapter source.\par\medskip
\noindent\textit{Allowed standard prerequisites (editorial).}\begin{itemize}
\item Branches over separably closed residue fields. More on Algebra, Definition more-algebra-definition-number-of-branches. Identifies the number of branches with the number of geometric branches when the residue field is separably closed.
\item The delta-branches inequality and its sharp case. Varieties, Lemmas varieties-lemma-delta-number-branches-inequality and varieties-lemma-delta-number-branches-inequality-sh. Supplies the general inequality and, in the strictly henselian equality case, the wedge-of-DVRs structure.
\item Invariance under completion and the coefficient field. More on Algebra, Lemma more-algebra-lemma-one-dimensional-number-of-branches; Varieties, Lemma varieties-lemma-delta-same-after-completion; Algebra, Lemma algebra-lemma-complete-henselian and algebra-lemma-regular-complete-containing-coefficient-field. Lets one pass to the completion without changing the invariants, get strict henselianness, and write each complete DVR branch as k[[t]].
\end{itemize}\medskip
\ifdefined\hypertarget\hypertarget{sourcepointer-1}{}\fi
\noindent\textit{Source reference.} Original lemma 04GM: sources/stacks-chow/tags.txt, source lines 5838--5838. Prerequisite pointer invoked at proof line(s) 3454; it names a standard fact of a sibling chapter that the counted argument uses, and every step of the counted proof built on it is retained verbatim in the delivered body.\par\medskip
\ifdefined\hypertarget\hypertarget{sourcepointer-2}{}\fi
\noindent\textit{Source reference.} Original lemma 0C0S: sources/stacks-chow/tags.txt, source lines 15097--15097. Prerequisite pointer invoked at proof line(s) 3458; it names a standard fact of a sibling chapter that the counted argument uses, and every step of the counted proof built on it is retained verbatim in the delivered body.\par\medskip
\ifdefined\hypertarget\hypertarget{sourcepointer-3}{}\fi
\noindent\textit{Source reference.} Original equation 0C1U: sources/stacks-curves/pinned-original-source.tex, source lines 3410--3410. Prerequisite pointer invoked at proof line(s) 3437; it names a standard fact of the same chapter that the counted argument uses, and every step of the counted proof built on it is retained verbatim in the delivered body.\par\medskip
\ifdefined\hypertarget\hypertarget{sourcepointer-4}{}\fi
\noindent\textit{Source reference.} Original definition 0C26: sources/stacks-chow/tags.txt, source lines 15146--15146. Prerequisite pointer invoked at proof line(s) 3444; it names a standard fact of a sibling chapter that the counted argument uses, and every step of the counted proof built on it is retained verbatim in the delivered body.\par\medskip
\ifdefined\hypertarget\hypertarget{sourcepointer-5}{}\fi
\noindent\textit{Source reference.} Original lemma 0C2D: sources/stacks-chow/tags.txt, source lines 15153--15153. Prerequisite pointer invoked at proof line(s) 3451; it names a standard fact of a sibling chapter that the counted argument uses, and every step of the counted proof built on it is retained verbatim in the delivered body.\par\medskip
\ifdefined\hypertarget\hypertarget{sourcepointer-6}{}\fi
\noindent\textit{Source reference.} Original lemma 0C43: sources/stacks-chow/tags.txt, source lines 15213--15213. Prerequisite pointer invoked at proof line(s) 3446; it names a standard fact of a sibling chapter that the counted argument uses, and every step of the counted proof built on it is retained verbatim in the delivered body.\par\medskip
\ifdefined\hypertarget\hypertarget{sourcepointer-7}{}\fi
\noindent\textit{Source reference.} Original lemma 0C42: sources/stacks-chow/tags.txt, source lines 15212--15212. Prerequisite pointer invoked at proof line(s) 3455; it names a standard fact of a sibling chapter that the counted argument uses, and every step of the counted proof built on it is retained verbatim in the delivered body.\par\medskip
\ifdefined\hypertarget\hypertarget{sourcepointer-8}{}\fi
\noindent\textit{Source reference.} Original lemma 0C3W: sources/stacks-chow/tags.txt, source lines 15206--15206. Prerequisite pointer invoked at proof line(s) 3452; it names a standard fact of a sibling chapter that the counted argument uses, and every step of the counted proof built on it is retained verbatim in the delivered body.\par\medskip
\begin{lemma}
\label{lemma-multicross-algebra}
Let $k$ be a separably closed field. Let $A$ be a $1$-dimensional
reduced Nagata local $k$-algebra with residue field $k$. Then
$$
\delta\text{-invariant }A \geq \text{number of branches of }A - 1
$$
If equality holds, then $A^\wedge$ is as in (\ref{equation-multicross}).
\end{lemma}

\begin{proof}
Since the residue field of $A$ is separably closed, the number
of branches of $A$ is equal to the number of geometric branches
of $A$, see
More on Algebra, Definition \ref{more-algebra-definition-number-of-branches}.
The inequality holds by
Varieties, Lemma \ref{varieties-lemma-delta-number-branches-inequality}.
Assume equality holds.
We may replace $A$ by the completion of $A$; this does
not change the number of branches or the $\delta$-invariant, see
More on Algebra, Lemma
\ref{more-algebra-lemma-one-dimensional-number-of-branches}
and Varieties, Lemma \ref{varieties-lemma-delta-same-after-completion}.
Then $A$ is strictly henselian, see
Algebra, Lemma \ref{algebra-lemma-complete-henselian}.
By Varieties, Lemma \ref{varieties-lemma-delta-number-branches-inequality-sh}
we see that $A$ is a wedge of complete discrete valuation rings.
Each of these is isomorphic to $k[[t]]$ by Algebra, Lemma
\ref{algebra-lemma-regular-complete-containing-coefficient-field}.
Hence $A$ is as in (\ref{equation-multicross}).
\end{proof}

\end{document}
